
\chapter{Prototypes} 

\ifpdf
    \graphicspath{{3_prototypes/figures/PNG/}{3_prototypes/figures/PDF/}{3_prototypes/figures/}}
\else
    \graphicspath{{3_prototypes/figures/EPS/}{3_prototypes/figures/}}
\fi



{\color{red}
\section{TODO}
add github and statistics
as well as explanations of github structure
}

\section{DQN Atari Breakout}
    
    For the DQN prototype, I implemented a Deep Q-Network agent to solve the Atari \textit{Breakout} environment. This implementation follows the core ideas presented by DeepMind and uses the \texttt{stable-baselines3} library to ensure a reliable and standardized implementation.



\subsection{Environment Setup}
    
    The Atari environment is created using the \textit{NoFrameskip} version of Breakout to maintain consistency with the original DQN experiments. Multiple parallel environments are instantiated to improve data collection efficiency. Frame stacking is applied so that the agent can observe temporal information across consecutive frames.
    
 \begin{lstlisting}[
    language=Python,
    numbers=left,
    numberstyle=\tiny,
    stepnumber=1,
    numbersep=5pt,
    breaklines=true,
    frame=shadowbox,
    basicstyle=\tiny\ttfamily,
    keywordstyle=\color{c++purple}\bfseries,
    stringstyle=\color{c++red},
    commentstyle=\color{c++green}
]
import ale_py
from stable_baselines3 import DQN
from stable_baselines3.common.env_util import make_atari_env
from stable_baselines3.common.vec_env import VecFrameStack
from stable_baselines3.common.callbacks import CheckpointCallback, EvalCallback


def main():
    vec_env = make_atari_env(
        "BreakoutNoFrameskip-v4",
        n_envs=8,
        seed=42,
    )

    vec_env = VecFrameStack(vec_env, n_stack=4)

    eval_env = make_atari_env(
        "BreakoutNoFrameskip-v4",
        n_envs=1,
        seed=123,
    )

    eval_env = VecFrameStack(eval_env, n_stack=4)
\end{lstlisting}

    

\subsection{Callbacks and Evaluation}
    
    Two callbacks are used during training. The checkpoint callback periodically saves the model parameters and replay buffer, enabling training to be resumed if needed. The evaluation callback assesses the agent’s performance at fixed intervals using a separate environment and stores the best-performing model based on evaluation rewards.
    
 \begin{lstlisting}[
    language=Python,
    numbers=left,
    numberstyle=\tiny,
    stepnumber=1,
    numbersep=5pt,
    breaklines=true,
    frame=shadowbox,
    basicstyle=\tiny\ttfamily,
    keywordstyle=\color{c++purple}\bfseries,
    stringstyle=\color{c++red},
    commentstyle=\color{c++green}
]
    checkpoint_callback = CheckpointCallback(
    save_freq=500_000,
    save_path="./checkpoints/",
    name_prefix="dqn_breakout",
    save_replay_buffer=True,
    save_vecnormalize=True,
    )
    
    eval_callback = EvalCallback(
    eval_env,
    best_model_save_path="./best_model/",
    log_path="./logs/",
    eval_freq=100_000,
    deterministic=True,
    render=False)
\end{lstlisting}
    

\subsection{DQN Model Configuration}
    
    The DQN agent uses a convolutional neural network policy suitable for image-based input. Hyperparameters such as learning rate, replay buffer size, discount factor, and exploration schedule are selected to closely match the settings used in the original DQN paper. A target network is updated at fixed intervals to stabilize learning, and the model is trained on a GPU to improve performance.
    
     \begin{lstlisting}[
    language=Python,
    numbers=left,
    numberstyle=\tiny,
    stepnumber=1,
    numbersep=5pt,
    breaklines=true,
    frame=shadowbox,
    basicstyle=\tiny\ttfamily,
    keywordstyle=\color{c++purple}\bfseries,
    stringstyle=\color{c++red},
    commentstyle=\color{c++green}
]
    model = DQN(
    "CnnPolicy",
    vec_env,
    learning_rate=1e-4,
    buffer_size=400_000,
    learning_starts=50_000,
    batch_size=64,
    tau=1.0,
    gamma=0.99,
    train_freq=4,
    gradient_steps=1,
    target_update_interval=10_000,
    exploration_fraction=0.1,
    exploration_initial_eps=1.0,
    exploration_final_eps=0.01,
    verbose=1,
    tensorboard_log="./tensorboard_logs/",
    device="cuda",
    )
    \end{lstlisting}
    

\subsection{Training Procedure}
    
    The agent is trained for a total of ten million timesteps. During training, both checkpointing and evaluation callbacks are active. After training is completed, the final model is saved and all environments are properly closed.
    
  \begin{lstlisting}[
    language=Python,
    numbers=left,
    numberstyle=\tiny,
    stepnumber=1,
    numbersep=5pt,
    breaklines=true,
    frame=shadowbox,
    basicstyle=\tiny\ttfamily,
    keywordstyle=\color{c++purple}\bfseries,
    stringstyle=\color{c++red},
    commentstyle=\color{c++green}
]
    total_timesteps = 10_000_000
    
    model.learn(
    total_timesteps=total_timesteps,
    callback=[checkpoint_callback, eval_callback],
    progress_bar=True,
    )
    
    model.save("dqn_breakout_final")
    vec_env.close()
    eval_env.close()
    
    
    if __name__ == "__main__":
    main()
    \end{lstlisting}
    

\subsection{Results}

    \figuremacroW{training_curves.png}{DQN Training Curves}{Training reward (top-left), evaluation reward (top-right), episode length (bottom-left), and loss (bottom-right) across training steps.}{0.7}

    % \figuremacroW{learning_curve_blob.png}{Title of Figure}{Description.}{0.7}
    
    \figuremacroW{eval_perfomance.png}{DQN Evaluation of Performance}{}{0.7} 
    
    \figuremacroW{eval_episode_length.png}{DQN Evaluation of Episode Length}{}{0.7} 
    

\subsection{Video / Images}

    link to YOUTUBE:
    \begin{itemize}
        \item \url{https://youtube.com/shorts/hdiVVfBSX9c?feature=share}
    \end{itemize}

    \figuremacroW{game_start.png}{Initial State of BreakoutNoFrameskip-v4}{Game state before the agent has taken any action.}{0.3} 
    \figuremacroW{game_end.png}{Game State Near End of a Trained Episode}{Environment state after sustained play, showing cleared bricks and accumulated score.}{0.3} 
